% =====================================================================
\section{G2 --- the gap, the region coverage, and the anatomy of the Savvidy cell}
\label{ch:gap}
% =====================================================================

\subsection{What G2 asks}

G2 asks for $f_1^2>0$ with isolation, equivalently $G''(0)\ge\lambda$ at $k\to0$, uniform in
$(L,a)$. The decomposition follows the corpus's assembly
(\code{research/unconditional/G2-ASSEMBLE-100.md}), with the labels of this paper.

\begin{table}[t]
\centering
\footnotesize
\begin{tabularx}{\textwidth}{@{}l>{\raggedright\arraybackslash}Xc>{\raggedright\arraybackslash}p{3.2cm}@{}}
\toprule
\textbf{Node} & \textbf{Object} & \textbf{Local \%} & \textbf{Label}\\
\midrule
G2.0 & Identity layer: $G''(k\to0)=E[H|0]-\operatorname{Var}(\nabla S|0)=V_{\mathrm{eff}}''=m_{\mathrm{eff}}^2(k\to0)$;
  exact Ward $E[a\,f]=1$; Stein/IBP $E[H]=E[f^2]$ and $E[\operatorname{Hess}V]=E[\nabla V\otimes\nabla V]\succeq0$ for every measure; $E[-D^2]=E[f^2]$ only for the free field; $G''\operatorname{Var}(\xi)=1$ iff the marginal is Gaussian; the sign sits in the \emph{conditioned} fibre Hessian & 100\% & \LT{} identities\\
G2.1 & Haar geometric term $W_{\mathrm{geo}}=E[\Phi'']-\operatorname{Var}(\Phi')\ge0$
  (Pr\'ekopa), every mode including $k\to0$, every $L$ & 100\% & \LT\\
G2.2 & Shallow-field band $\{2g\|F\|\le N/12\}$ closed by $\Phi''\succeq N/12$
  $+$ Brascamp--Lieb & 100\% & \LT\\
G2.3 & High-field cone (WSTC): cone variance mass
  $\le(c\beta\kappa_{\max})^2e^{-c\beta(\kappa^2-\kappa_{\rm typ}^2)N_\square}\to0$,
  $L$-uniform & 100\% & \LT\\
G2.4 & AF-UV deep cone: instability $\le\mu^2e^{-c\mu^4/g^4}\to0$ & 100\% & \LT\\
\textbf{G2.5} & \textbf{Open cell $=$ typical field $\times$ high momentum}:
  the ceiling $\operatorname{Var}(\xi_{k\to0})\le1/\lambda$ (curvature form
  $\operatorname{Var}_\mu(\langle v,\nabla S\rangle|0)\le E_\mu[H_{vv}|0]+W_{\mathrm{full}}-\delta$) &
  \textbf{0\%} & \textbf{\LM{} $=$ H1}\\
\textbf{G2.6} & \textbf{Named species}: negative association / anti-FKG
  $\operatorname{Cov}(G,\Phi'|0)\le0$ & \textbf{0\%} &
  \textbf{\LM{} proven absent from nine classes}\\
\bottomrule
\end{tabularx}
\caption{G2 node board: the gap.}
\label{tab:g2}
\end{table}

The region coverage of the $(\text{field}\times\text{momentum})$ plane is about \textbf{90\%},
closed by the theorems G2.1--G2.4 (Figure~\ref{fig:plane}). The remaining one cell --- typical
field, high momentum --- is a region where all four handles fail \emph{simultaneously}: the WSTC
exponent vanishes when the field is typical, the AF-UV suppression has not begun at
$g\approx1.5$, and the $N/12$ floor is roughly ten times too small. That cell is
\textbf{exactly} $\operatorname{Var}(\xi_{k\to0})\le1/\lambda$, with zero slack. The sign --- the
actual proof --- is \textbf{0\%}.

The G2 assembly is complete in the following precise sense: of 26 input rows, every row is
either (a) a theorem or (b) a reduction to the single species; \textbf{there is no row (c)}, a
genuinely independent open item. The curvature face H1 ($G''\ge\lambda$, the \emph{inverse
zero-momentum susceptibility} $1/\chi$) and the spectral-edge face $H^*$ ($f_1^2\ge r_0$ with
isolation I2) are \emph{not} one object: by the squeeze $f_1^2/\chi\le m_1^2\le M_0/\chi$, H1 is a
susceptibility that \emph{upper}-bounds $m_1^2$ and is necessary but strictly weaker. The
identification $H^*=\text{H1}$ is not merely unproved but is \emph{contradicted} for composite
operators --- the free massive $\mathord{:}\varphi^2\mathord{:}$ has a gap ($m_1=2m$) yet
$G''=1/\chi\ne m_1^2$ --- and every gauge-invariant Yang--Mills operator is composite. The gap is
$H^*$; the rigorous content is the residue--level-gap lemma and the exact Ward identity $E[a\,f]=1$.

\figplane

\subsection{The open cell, stated exactly}
\label{sec:opencell}

It is useful to write the open cell three times, because the three forms are used in different
chapters and the reader should see that they are one inequality:
\begin{align}
\text{(variance form)}\qquad & \operatorname{Var}_\mu\bigl(\xi_{k\to0}\bigr)\ \le\ 1/\lambda ;
\label{eq:opencell-var}\\
\text{(curvature form)}\qquad &
  \operatorname{Var}_\mu\bigl(\langle v,\nabla S\rangle\mid0\bigr)
  \ \le\ E_\mu[H_{vv}\mid0]+W_{\mathrm{full}}-\delta ;
\label{eq:opencell-curv}\\
\text{(defect form)}\qquad & D\ \le\ S-\lambda\quad\text{together with}\quad S\ \ge\ \lambda ,
\label{eq:opencell-defect}
\end{align}
where $S$ and $D$ are the exact variance-decomposition terms of Chapter~\ref{sec:theorems}. The
first is H1; the second is its invariant (conditioned) form; the third is the form in which the
Savvidy cell is discussed below. The bare defect form $D\le(1-\epsilon)S$ is \emph{insufficient}:
it allows $S$ to be arbitrarily small, and the gap needs the absolute floor $S\ge\lambda$.

\textbf{Why the cell is where all handles fail.} The four closed cells are closed by four
different mechanisms: low field by the uniform Haar curvature $\Phi''\succeq N/12$ plus
Brascamp--Lieb; high field and low momentum by measure starvation (the WSTC exponent grows with
the field excess); high field and high momentum by asymptotic freedom (the instability is
$e^{-c\mu^4/g^4}$); and the AF-UV deep cone by the same exponent. In the remaining cell ---
typical field, coherence length near $k\to0$ but momentum high enough that the field is not
exponentially suppressed --- every exponent vanishes at once. This is not a gap in the
argument; it is the exact intersection of the four mechanisms' domains of validity, and it is
nonempty precisely at the physical coupling.

\subsection{The negative-association species}
\label{sec:negspecies}

The variance-free face (the only one that could pierce H1) is a deterministic
\textbf{negative-association / anti-FKG} inequality
\begin{equation}
\operatorname{Cov}_\mu\bigl(G,\Phi'\mid0\bigr)\ \le\ 0
\label{eq:neqassoc}
\end{equation}
for the continuous, frustrated (cos-sum Wilson action), non-attractive $\mathrm{SU}(N)$
Haar $\times$ Wilson measure. The campaign proves this object is \textbf{absent from every known
dependence class} (\code{C5-NEG-ASSOC.md}, \code{C6-DEPENDENCE-CLASS.md}): negative association
(NA), negative regression/NOD, strong Rayleigh, log-submodular, determinantal (DPP), total
positivity (TP$_2$; Chebyshev--Fr\'echet), FKG/Griffiths, Brascamp--Lieb, and Simon--Lieb all
fail, each for a specific structural reason. The reasons are worth listing, because they are the
campaign's clearest catalogue of why ``known positive dependence'' does not apply:

\begin{itemize}
\item \textbf{No order on the quotient.} Negative association and FKG both need a lattice order;
  the gauge quotient $\mathcal{A}/\mathcal{G}$ carries no order, so the definitions do not even
  parse (this is the nested-cone no-go of Chapter~\ref{sec:nogos}).
\item \textbf{No real-stable polynomial.} Strong Rayleigh and determinantal structures require a
  real-stable generating polynomial or a DPP kernel; the $\mathrm{SU}(N)$ Weingarten weights are
  signed, so the polynomial is not real-stable.
\item \textbf{No common scalar.} Log-submodularity needs a common scalar order; the matrix-valued
  field has no such order.
\item \textbf{FKG wrong-signed.} The Wilson action is ferromagnetic in the plaquette variable but
  the observable is a cos-sum of traces; the monotone coupling is not the one that would give
  positive association of $G$ and $\Phi'$.
\item \textbf{Brascamp--Lieb needs a convex potential}, which is false at Savvidy
  (Chapter~\ref{sec:savvidy}); and Simon--Lieb is sign-blind on the variance, which is exactly
  H1; the Stein--Hellfer--Sj\"ostrand constant is $1/\lambda$, again H1.
\end{itemize}

The measured $\operatorname{Cov}\approx-0.11$ at $\beta=2.5$ is the chart-local log-cut shadow of
\emph{proven} $\Phi$-convexity, not a law: the per-mode $X\in[-0.09,+0.16]$ changes sign, and
$\beta=5.7$ flips it. So the species is not merely unproven; its sign is not stable mode by mode,
which is consistent with its being a genuine dynamical statement rather than a kinematic
identity.

\subsection{Anatomy of the Savvidy cell}
\label{sec:savvidy}

The Savvidy cell is the one region where the classical one-loop effective potential has an
indefinite Hessian, and for two decades it was treated in the campaign's own notes as the
``true physical obstruction.'' This edition records the wave that dismantled that reading. The
outcome is a genuine theorem-set that \emph{removes a false obstruction}, and a precise
statement of what survives. We state the removal first, then the residue.

\subsubsection{The false obstruction, and why it is false}

In a constant background $F_B$, the one-loop effective action has a negative mode: the Hessian
$-D^2+2\operatorname{ad}F_B$ has a negative eigenvalue in the transverse-magnetic sector. The
naive reading is that the Savvidy vacuum is unstable, that the effective action is concave, and
that any variance or curvature bound on the coarse mode must fail. The wave
(\code{SAVVIDY-DTERM.md}, \code{SAVVIDY-NEGMODE-EXACT.md}, \code{SAVVIDY-VAR-PHI2.md}) shows
that this reading is wrong in three separate senses.

\textbf{(i) The negative mode is a barrier-top artifact.} The one-loop potential is a double
well with the true vacuum at $\varphi=\pm v$; the negative mode sits at the barrier top
$\varphi=0$, not at the vacuum. At the true vacuum the curvature is positive, $V''=2c>0$, and
the fluctuation spectrum is stable. The indefinite Hessian is a property of a saddle, not of the
measure. This is a \LT{} statement about the effective potential, and it is the reason the
periodic no-go (Chapter~\ref{sec:nogos}) does not apply to the true vacuum.

\textbf{(ii) The softening is bounded and $c$-independent.} The fluctuation of the order
parameter is
\begin{equation}
\operatorname{Var}(\varphi^2)\;=\;\frac{2}{\lambda_4}+\frac{2}{(gB)^2}+\cdots
\;\longrightarrow\;\frac{2}{\lambda_4},
\label{eq:varphi2}
\end{equation}
bounded and independent of the vacuum parameter $c$. The softened curvature is
$\Gamma''(0)=\kappa-g^2\operatorname{Var}(\varphi^2)$, and the softening term $g^2
\operatorname{Var}(\varphi^2)$ is $O(1)$ and $g$-independent: it does not diverge as $g\to0$,
and it does not grow with the background. The earlier fear that the negative mode produces a
divergence was a misreading of the one-loop determinant.

\textbf{(iii) The sign is protected at two loops.} The variance is a positive quadratic form
times $g^2$; it renormalizes multiplicatively; the two-loop correction moves the constant by
$\lesssim1\%$ and never crosses zero (\code{SAVVIDY-2LOOP.md}). So the softening cannot change
sign under the renormalization group within the perturbative regime.

\subsubsection{The exact $D$-split}

The variance decomposition $G''=S-D$ of Chapter~\ref{sec:theorems} has, in the Savvidy cell, an
exact split (\code{SAVVIDY-DTERM.md}):
\begin{equation}
D\;=\;D_+(\le0)\;+\;\frac{\mathbb E\bigl[|\langle\varphi_-|c\rangle|^2\bigr]}{g|B|},
\label{eq:dsplit}
\end{equation}
where $\varphi_-$ is the rank-one negative mode and $D_+\le0$ is the positive-semidefinite part.
The inequality $D<S$ --- the open cell --- is therefore equivalent to the rank-one term being
smaller than the bare stiffness $S_{\mathrm{bare}}$. In the deep infrared the rank-one term
decays and $D/S\to0$; at the crossover it is $O(1)$. This is the sharpest form of the Savvidy
statement in the campaign: \emph{the obstruction is rank-one}, and the question is whether that
one direction can be dominated uniformly.

\subsubsection{Uniformity, and the tension with the full flow}

Two results pull in opposite directions, and the paper records both rather than choosing.
\begin{itemize}
\item \emph{One-octave uniformity} (\code{SAVVIDY-UNIFORM-SOFTENING.md}): the soft modes have
  $|p|\sim\sqrt{gB}$, a single octave; there is no logarithmic divergence across octaves; hence
  the softening is $(L,a)$-uniform. On this reading the Savvidy cell contributes one marginal
  $O(1)$ inequality at the crossover.
\item \emph{Multi-cell accumulation} (\code{SAVVIDY-MULTICELL-RG.md}): integrating the softening
  over the full flow accumulates a $\ln\ln$ growth near the Landau pole, which overshoots for
  $O(1)$ bare stiffness. On this reading the cell is \emph{infrared-marginal}, and the question
  is decided nonperturbatively at $k_{\mathrm{IR}}\sim\Lambda$, which is H1.
\end{itemize}
The reconciliation is that the per-octave softening is genuinely bounded (a theorem), while the
full-flow accumulation is infrared-dominated (also a theorem). So the Savvidy cell is not ``one
easy $O(1)$ inequality''; it is a nonperturbative infrared statement wearing the costume of a
one-loop computation. This is a genuine sharpening of the RG missing step, not a closure.

\subsubsection{The lattice cell}

The lattice measurement (\code{SAVVIDY-LATTICE-CELL.md}) gives the numbers that make the
near-criticality concrete. The defect ratio $D/S$ is bounded ($\le1.1$) but $O(1)$:
$0.94\le D/S\le1.00$. Five of twelve modes have $G''<0$. The $k\to0$ mode is absent on the
available volumes. The earlier reported $G''\approx+0.28$ is a small-volume artifact of the
free dispersion $4\sin^2(\pi/L)$: on a small torus the lowest momentum is not small enough to
be a continuum probe. So the lattice does not contradict the theorem-set, and it does not
confirm a positive margin either. It shows the margin is thin.

\figsavcell

\subsubsection{What is removed and what remains}

\begin{itemize}
\item \textbf{Removed (theorems):} the claim that the Savvidy negative mode causes a divergence
  or destroys the $\chi$ floor. The negative mode is a barrier-top artifact; the true vacuum is
  stable; the softening is bounded by $2/\lambda_4$, $c$-independent, sign-protected at two
  loops, and one-octave uniform; and the deep-infrared $D/S\to0$.
\item \textbf{Not removed (still H1):} the surviving sub-criticality is infrared-marginal in the
  full RG flow and empirically near-critical. Establishing a positive uniform margin is exactly
  the open cell, written in the Savvidy coordinate.
\end{itemize}

This is the campaign's cleanest example of a \emph{false} obstruction being removed while the
true wall stays put. It is also a warning about the paper's own earlier drafts: the
``$D/S\approx2\%$'' figure belongs to a gauge-dependent gluon mode and a weak-coupling regime, not
to the gauge-invariant $0^{++}$ object at the physical coupling (Appendix~\ref{app:postaudit}).

\subsection{The gauge quotient and the $7.1\%$ residue}

The negativity that survives the gauge quotient is measurable. After projecting out the gauge
directions ($Q_\perp$), the numerical experiment W4 finds $7.1\%$ negative directions in the
fluctuation spectrum at the physical coupling. This is the number the paper quotes when it says
the wall is real rather than a gauge artifact: the transverse-magnetic negative mode survives
the quotient. It is the same mode that the periodic no-go identifies as the price of winding on
the compact group, and the same mode whose squared overlap with the $k\to0$ coarse mode vanishes
as $(k/\sqrt{gB})^2$ --- which is why the coarse-mode softening is bounded even though the
negative mode exists.

\subsection{Summary}

Region coverage $\sim90\%$, by genuine theorems. The sign is 0\%, by one inequality, in one cell.
The Savvidy cell is where that cell is sharpest, and the wave summarized above shows both what
the cell is \emph{not} (a divergence, a gauge artifact, an easy one-loop inequality) and what it
\emph{is} (a rank-one, near-critical, infrared-marginal sub-criticality). The next chapter
collects the proven theorems that close the other 90\%.
